\chapter{Zasiegi, bloki i widocznosc nazw}

\section{Dlaczego zasieg jest tak wazny}
Zasieg odpowiada na pytanie: \emph{z jakiego miejsca programu dana nazwa jest widoczna
i do jakiej wartosci sie odnosi}. Bez tego bardzo szybko powstalby chaos. Dwie rozne
czesci programu moglyby przypadkiem walczyc o te sama nazwe i nadpisywac sobie dane.

W MoonChunk temat zasiegow jest szczegolnie istotny, bo laczy:

\begin{itemize}
  \item zwykle zmienne,
  \item funkcje,
  \item ciała warunkow i petli,
  \item odrebne chunki,
  \item jawne bloki \mc{\{ ... \}},
  \item specjalny zapis \mc{parent::name}.
\end{itemize}

\section{Najprostsza intuicja}
Jesli deklarujesz nazwe w pewnym miejscu, to najpierw nalezy zakladac, ze jest ona
widoczna lokalnie w tym miejscu i w tym, co z niego wynika. Jednak nie kazdy blok w
MoonChunk daje identyczna swobode cieniowania nazw.

To celowy projekt: jezyk stawia na przewidywalnosc i chce ograniczac przypadkowe
kolizje nazw.

\section{Redeclaracja a zwykle przypisanie}
Na poczatek trzeba odroznic dwie rzeczy:

\begin{itemize}
  \item \important{redeklaracja} tworzy nowa nazwe o tej samej nazwie,
  \item \important{przypisanie} zmienia wartosc juz istniejacej nazwy typu \mc{let}.
\end{itemize}

Przyklad poprawnego przypisania:

\begin{lstlisting}[style=moonchunk,caption={Zmiana wartosci przez przypisanie}]
let x: int = 1;
x = x + 1;
\end{lstlisting}

Przyklad błednej redeklaracji w tym samym zakresie:

\begin{lstlisting}[style=moonchunk,caption={Bledna redeklaracja w tym samym miejscu}]
let x: int = 1;
let x: int = 2;
\end{lstlisting}

\section{Granice, przy ktorych ta sama nazwa jest dozwolona}
MoonChunk dopuszcza powtorzenie tej samej nazwy tylko w wybranych, swiadomych granicach.
Najwazniejsze z nich to:

\begin{itemize}
  \item osobny \mc{chunk},
  \item osobne wywolanie funkcji,
  \item jawny blok \mc{\{ ... \}}.
\end{itemize}

Oznacza to, ze ten sam identyfikator moze oznaczac co innego w innym chunku albo wewnatrz
funkcji, bez konfliktu z nazwa z zewnatrz.

\section{Osobne chunki jako osobne granice}
Dwa rozne chunki moga korzystac z tej samej nazwy i nie wchodza sobie w droge.

\begin{lstlisting}[style=moonchunk,caption={Ta sama nazwa w dwoch chunkach}]
chunk "First" {
  const x: int = 1;
  print("First.x:", x);
};

chunk "Second" {
  const x: int = 2;
  print("Second.x:", x);
};
\end{lstlisting}

To naturalne: kazdy chunk jest osobna jednostka wykonania.

\section{Funkcja jako granica zasiegu}
Funkcje rowniez tworza istotna granice. Przyklad:

\begin{lstlisting}[style=moonchunk,caption={Ta sama nazwa na zewnatrz i wewnatrz funkcji}]
chunk "Main" {
  const x: int = 3;

  function demo(): int {
    const x: int = 5;
    return x;
  }

  print("outer:", x);
  print("inner:", demo());
};
\end{lstlisting}

Tutaj wewnetrzne \mc{x} funkcji nie niszczy zewnetrznego \mc{x}. To oddzielna przestrzen.
Dodatkowo przy rekurencji kazde kolejne wywolanie funkcji ma swoj wlasny zestaw lokalnych
zmiennych.

\section{Jawny blok \mc{\{ ... \}}}
MoonChunk wspiera tez samodzielny, jawny blok, ktory pozwala swiadomie utworzyc nowy
zagniezdzony zakres.

\begin{lstlisting}[style=moonchunk,caption={Jawny blok jako osobny zakres}]
let i: int = 1;

{
  let i: int = 99;
  print(i);
};

print(i);
\end{lstlisting}

W takim przypadku wewnetrzne \mc{i} jest osobna nazwa dla wnetrza bloku, a zewnetrzne
\mc{i} pozostaje bez zmian.

\section{Dlaczego to rozwiazanie jest przydatne}
Jawny blok pozwala:

\begin{itemize}
  \item ograniczyc zycie zmiennych do małego fragmentu logiki,
  \item eksperymentowac z nazwami pomocniczymi bez zanieczyszczania otoczenia,
  \item budowac czytelniejsze etapy obliczen,
  \item tworzyc bezpieczne, lokalne konteksty pracy.
\end{itemize}

To bardzo przyjemny mechanizm wszedzie tam, gdzie funkcja bylaby zbyt duza konstrukcja,
a zwykly ciag instrukcji -- zbyt splaszczony.

\section{Ciala \mc{if}, \mc{for} i \mc{while} sa bardziej rygorystyczne}
To wazna cecha MoonChunk. Chociaz instrukcje warunkowe i petle korzystaja z blokow,
nie nalezy traktowac ich jak pelnej, swobodnej przestrzeni do ponownego deklarowania
tej samej nazwy z najblizszego otoczenia.

Przyklad problematyczny:

\begin{lstlisting}[style=moonchunk,caption={Redeclaracja w ciele if}]
const x: int = 3;

if (true) {
  const x: int = 5;
};
\end{lstlisting}

Ten wzorzec powinien byc traktowany jako blad redeklaracji. Jesli celem jest swiadome
wydzielenie osobnego zasiegu z ta sama nazwa, lepiej uzyc jawnego bloku \mc{\{ ... \}}
albo funkcji.

\section{Odczyt wartosci z nadrzednego zakresu przez \mc{parent::}}
MoonChunk daje tez specjalny zapis pozwalajacy siegnac do wyzszej warstwy zasiegu.
Przyklad:

\begin{lstlisting}[style=moonchunk,caption={Odczyt z parent::}]
let i: int = 1;

{
  let i: int = 99;
  {
    let fromParent1: int = parent::i;
    let fromParent2: int = parent::parent::i;
    print(fromParent1);
    print(fromParent2);
  };
};
\end{lstlisting}

Tutaj:

\begin{itemize}
  \item \mc{parent::i} oznacza \mc{i} z bezposrednio wyzszego poziomu,
  \item \mc{parent::parent::i} oznacza \mc{i} dwa poziomy wyzej.
\end{itemize}

To narzedzie zaawansowane, ale bardzo przydatne, gdy swiadomie pracujesz na zagniezdzonych
blokach.

\section{Kiedy nie uzywac \mc{parent::}}
Jesli mozna przekazac wartosc wprost do funkcji albo zbudowac kod czytelniej przez
zmienne pomocnicze, zwykle warto wybrac prostsze rozwiazanie. \mc{parent::} jest
przydatny, ale naduzywany szybko utrudnia sledzenie przeplywu danych.

\section{Globalne wartosci a zasiegi lokalne}
Zmienne globalne z bloku \mc{env} sa osobnym mechanizmem. Sluza do wspolnych wartosci,
ale nie powinny byc traktowane jako substytut dobrej organizacji lokalnych nazw.

Praktyczna rada jest prosta:

\begin{itemize}
  \item rzeczy naprawde wspolne wrzucaj do globali,
  \item rzeczy lokalne zostaw lokalnie,
  \item nie buduj wszystkiego na globalach tylko dlatego, ze sa wygodne.
\end{itemize}

\section{Najczestsze bledy zwiazane z zasiegami}
\begin{itemize}
  \item redeklaracja tej samej nazwy w tym samym miejscu,
  \item probowanie cieniowania nazwy w \mc{if}, \mc{for} lub \mc{while},
  \item odczyt zmiennej przed inicjalizacja,
  \item zbyt glebokie poleganie na \mc{parent::},
  \item mieszanie odpowiedzialnosci globali i zwyklych zmiennych lokalnych.
\end{itemize}

\section{Zlota zasada czytelnosci}
Jesli po kilku dniach nie jestes w stanie szybko odpowiedziec, skad bierze sie wartosc
jakiejs nazwy, to znaczy, ze struktura zasiegow jest zbyt skomplikowana. Dobrze napisany
kod MoonChunk ma byc czytelny nie tylko dla interpretera, ale tez dla czlowieka.
